% SGA 3，第 VI_B 篇，第一轮英文重建稿的中文翻译。
% 范围：推论10.11.1、命题10.12及其完整证明、推论10.12.1、
% 推论10.13及其完整证明，含编者注98--101。
% 依据：Polo--Gille Exp6B-13oct24.pdf，本组件局部页73--74
% （合卷阅读版页414--415；再版页眉403--404；原版页边印刷页392--394）。
% 段落10.14始于局部页74，不在本组件内。
% OCR及Reinhold英文Markdown仅用作对照见证。
% 本限界组件不含实际图。

\medskip
\phantomsection\label{cor:sga3-VIB-10.11.1}%
\noindent\textbf{推论 10.11.1.---}
\begingroup\itshape
设 \(S\) 是泛点为 \(\eta\) 的整概形。
设 \(G,H,K\) 是有限呈示的 \(S\)-群，
\[
  i:H\longrightarrow G,
  \qquad
  j:K\longrightarrow G
\]
是两个拟紧的 \(S\)-群单态射。则存在 \(S\) 的非空开子集 \(S'\)，
使经基变换得到的态射
\[
  i':H'\longrightarrow G',
  \qquad
  j':K'\longrightarrow G',
\]
是闭浸入，并使函子
\phantomsection\label{form:sga3-VIB-10.11.1-transporter-functors}%
\[
  \underline{\operatorname{Transp}}_{G'}(H',K'),
  \qquad
  \underline{\operatorname{Transp}}^{\mathrm{str}}_{G'}(H',K'),
\]
分别地，函子
\phantomsection\label{form:sga3-VIB-10.11.1-centralizer-normalizer}%
\[
  \underline{\operatorname{Centr}}_{G'}H',
  \qquad
  \underline{\operatorname{Norm}}_{G'}H',
\]
由 \(G'\) 中在 \(S'\) 上有限呈示的闭 \(S'\)-子概形
\textup{（分别地，闭 \(S'\)-子群）}表示。
\par\endgroup

这由\hyperref[prop:sga3-VIB-10.11]{上一命题}推出，
因为由 \hyperref[cor:sga3-VIB-1.4.2]{1.4.2}，假设蕴含
\(i_\eta\) 和 \(j_\eta\) 是闭浸入。

\medskip
\phantomsection\label{prop:sga3-VIB-10.12}%
\noindent\textbf{命题 10.12.---}
\begingroup\itshape
设 \(S\) 是整概形，\(G\) 是有限呈示的 \(S\)-群，
\(A\) 和 \(B\) 是 \(G\) 的两个群子概形，在 \(S\) 上有限呈示，且泛纤维平滑。
再假设下列条件之一成立：
\begin{enumerate}
\item[\textup{a)}]\phantomsection\label{cond:sga3-VIB-10.12-a}
  \(A\) 的泛纤维连通；
\item[\textup{b)}]\phantomsection\label{cond:sga3-VIB-10.12-b}
  \(A\) 和 \(B\) 在 \(G\) 中正规。
\end{enumerate}
则存在 \(S\) 的非空开子集 \(S'\)，以及
\[
  G'=G|_{S'},
\]
的闭群子概形 \(D'\)，在 \(S'\) 上有限呈示且纤维平滑，
它表示由 \(G'\) 中
\[
  A'=A|_{S'}
  \qquad\text{和}\qquad
  B'=B|_{S'}
\]
的换位子群的预层所伴随的 \textup{(fppf)} 层。
此外，在情形\textup{(a)}中，\(D'\) 的纤维连通；
在情形\textup{(b)}中，它在 \(G'\) 中正规。

特别地，对每个 \(s\in S'\)，用 \hyperref[def:sga3-VIB-7.2.2]{7.2.2} 的记号，有
\[
  D'_s=(A_s,B_s),
\]
\par\endgroup

\emph{证明。}
设 \(\eta\) 是 \(S\) 的泛点，并令
\[
  H_\eta=(A_\eta,B_\eta).
\]
由于 \(A_\eta\) 和 \(B_\eta\) 平滑，在情形\textup{(a)}中由
\hyperref[prop:sga3-VIB-7.8]{7.8}，在情形\textup{(b)}中由
\hyperref[concl:sga3-VIB-7.3-v]{7.3\textup{(v)}}，可知
\(H_\eta\) 分别为连通的、在 \(G_\eta\) 中正规的。

这里处于 \hyperref[par:sga3-VIB-10.0]{10.0} 中对应于
\hyperref[par:sga3-VIB-10.10-a]{10.10\textup{(a)}} 的情形。
从而，由 \hyperref[lem:sga3-VIB-10.3]{10.3} 和
\hyperref[par:sga3-VIB-10.1]{10.1}，存在 \(S\) 的非空开子集
\(S'\) 及有限呈示、在 \(G'\) 中闭的 \(S'\)-群子概形 \(D'\)，使
\[
  D'_\eta=D'\otimes_{S'}\kappa(\eta)=H_\eta.
\]
此外，由 EGA~IV\(_3\)，9.7.7和9.3.3，可以假定 \(D'\) 的纤维几何约化。
在情形\textup{(a)}中，再次应用 EGA~IV\(_3\)，9.7.7和9.3.3，
可假定 \(D'\) 的纤维连通，因而几何连通
\textup{（参见 \hyperref[prop:sga3-VIA-2.1.1]%
{\SGARefLink{prop:sga3-VIA-2.1.1}{VI\(_A\)，2.1.1}}）}。
在情形\textup{(b)}中，由 \hyperref[lem:sga3-VIB-10.4]{10.4}，
可假定 \(D'\) 在 \(G'\) 中正规。

我们在 \hyperref[prop:sga3-VIB-7.8]{7.8} 的证明中还已看到，存在整数
\(n\)，使
\phantomsection\label{form:sga3-VIB-10.12-generic-image}%
\[
  \nu_\eta^n
  \left((A_\eta\times_{\kappa(\eta)}B_\eta)^n\right)
  =D'_\eta,
\]
其中 \(\nu_\eta\) 和 \(\nu_\eta^n\) 如
\hyperref[def:sga3-VIB-7.2.2-a]{7.2.2\textup{(a)}} 及
\hyperref[concl:sga3-VIB-7.1-ii]{7.1\textup{(ii)}} 中所定义。
相同的公式定义态射
\phantomsection\label{form:sga3-VIB-10.12-nu-prime}%
\[
  \nu':A'\times_{S'}B'\longrightarrow G',
  \qquad
  {\nu'}^n:(A'\times_{S'}B')^n\longrightarrow G',
\]
且
\[
  {\nu'}^n\otimes_{S'}\kappa(\eta)=\nu_\eta^n.
\]

由 \hyperref[par:sga3-VIB-10.1]{10.1}，可以选取 \(S'\)，使态射
\({\nu'}^n\) 平坦且通过 \(D'\) 分解，并使所得态射
\[
  {\nu'}^n:(A'\times_{S'}B')^n\longrightarrow D'
\]
为满射。由 \hyperref[lem:sga3-VIB-7.5]{7.5}，态射%
\footnote{\label{note:sga3-VIB-98}编者注：下文的 \(n+1\) 已更正为 \(2n\)。}
\phantomsection\label{form:sga3-VIB-10.12-covering-map}%
\[
  {\nu'}^{2n}
  =\mu\circ\bigl({\nu'}^n\times_{S'}{\nu'}^n\bigr):
  (A'\times_{S'}B')^{2n}\longrightarrow D'
\]
是 \textup{(fppf)} 拓扑的覆盖。因此，由
\hyperref[prop:sga3-VIB-7.6]{7.6}，\(D'\) 表示由 \(G'\) 中
\(A'\) 和 \(B'\) 的换位子预层所伴随的 \textup{(fppf)} 层。

进一步，对每个 \(s\in S'\)，\({\nu'}^n\) 诱导满射态射
\[
  \nu_s^n:
  (A_s\times_{\kappa(s)}B_s)^n\longrightarrow D'_s.
\]%
\footnote{\label{note:sga3-VIB-99}编者注：对原文接下来的论证作了扩充。}
因此 \(D'_s\) 是 \(G_s\) 的闭子群，包含
\[
  \nu_s(A_s\times_{\kappa(s)}B_s),
\]
因而包含 \((A_s,B_s)\)。由于 \(\nu_s^n\) 是满射，有
\(D'_s=(A_s,B_s)\)；又由 \hyperref[prop:sga3-VIB-7.6]{7.6}，
它表示 \(G_s\) 中 \(A_s\) 和 \(B_s\) 的换位子的 \textup{(fppf)} 层。

\medskip
\phantomsection\label{cor:sga3-VIB-10.12.1}%
\noindent\textbf{推论 10.12.1.---}%
\footnote{\label{note:sga3-VIB-100}编者注：添加了本推论，
它用于 \hyperref[prop:sga3-VIB-8.4]{8.4} 的证明。}
\begingroup\itshape
设 \(S\) 是泛点为 \(\eta\) 的整概形，\(G\) 是纤维平滑的有限呈示
\(S\)-群。对每个 \(p\in\mathbb N^*\)，令
\[
  K_\eta^0=G_\eta,
  \qquad
  K_\eta^p=(K_\eta^{p-1},K_\eta^{p-1})
  \quad
  \textup{（分别地，}
    K_\eta^p=(G_\eta,K_\eta^{p-1})\textup{）}
\]
。%
\begingroup
\renewcommand{\thefootnote}{\fnsymbol{footnote}}%
\footnote[1]{\label{note:sga3-VIB-10.12.1-Geta}
译者注：法文再版的第二个递推式印作 \(G\)。由于 \(K_\eta^{p-1}\)
是 \(G_\eta\) 的 \(\kappa(\eta)\)-群子概形，必须与 \(G_\eta\) 作换位子。}%
\endgroup
固定 \(n\in\mathbb N^*\)。则存在 \(S\) 的非空开子集 \(S'\)
及 \(G|_{S'}\) 中正规的群子概形 \(D\)，它有限呈示且纤维平滑，
使对每个 \(s\in S'\)，有
\[
  D_s=K_s^n
\]
\par\endgroup

这由 \hyperref[prop:sga3-VIB-10.12]{10.12} 对 \(n\) 作归纳得到。

\medskip
\phantomsection\label{cor:sga3-VIB-10.13}%
\noindent\textbf{推论 10.13.---}
\begingroup\itshape
设 \(S\) 是泛点为 \(\eta\) 的整概形。
设 \(G\) 是 \(S\)-群，\(H\) 是 \(G\) 的正规群子概形。
假设 \(G\) 和 \(H\) 在 \(S\) 上有限呈示，且泛纤维平滑。%
\footnote{\label{note:sga3-VIB-101}编者注：原文假设 \(G\) 和 \(H\)
有限呈示，且 \(H\) 的纤维平滑；为了应用
\hyperref[prop:sga3-VIB-10.12]{10.12}，修改了假设。证明也作了扩充。}

若
\[
  (H_\eta,H_\eta)=H_\eta
  \qquad
  \textup{（分别地，}
    (G_\eta,H_\eta)=H_\eta\textup{）},
\]
则存在 \(S\) 的非空开子集 \(S'\)，使对每个 \(s\in S'\)，有
\[
  (H_s,H_s)=H_s
  \qquad
  \textup{（分别地，}
    (G_s,H_s)=H_s\textup{）}.
\]
\par\endgroup

事实上，\hyperref[prop:sga3-VIB-10.12]{10.12} 的证明给出了
\(S\) 的非空开子集 \(S'\) 及 \(G|_{S'}\) 的 \(S'\)-群子概形
\(D\)，它有限呈示且纤维平滑，使对每个 \(s\in S'\)，有
\[
  D_s=(H_s,H_s)
  \qquad
  \textup{（分别地，}
    D_s=(G_s,H_s)\textup{）}.
\]
另一方面，由于 \(D_\eta=H_\eta\)，且 \(D\) 和 \(H\) 在 \(S'\)
上有限呈示，EGA~IV\(_3\)，8.8.2.5 给出 \(S'\) 的非空开子集
\(S''\)，使
\[
  D|_{S''}=H|_{S''}.
\]
从而，对每个 \(s\in S''\)，有
\[
  H_s=D_s=(H_s,H_s)
  \qquad
  \textup{（分别地，}
    H_s=D_s=(G_s,H_s)\textup{）}.
\]

% 中文翻译检查点：本组件全部正文、证明、脚注及源注已译；源287行。
% 机械核对：行内数学118处、陈列数学23处及全部引用键守恒（数学中的自然语言已译）。
% 状态：完成源文覆盖和逐段自校；尚未经独立合卷编译与页面审查。
